\clearpage

\appendix

\begin{figure}[t!]
    \centering
    \includegraphics[width=0.85\columnwidth]{Figures/discipline_fig.pdf}
    \vspace{-0.075in}
    \caption{Visualization of the distribution of disciplines for all core papers, selected for research idea generation.}
    \label{fig:discipline}
    \vspace{-0.1in}
\end{figure}

\section{Additional Experimental Details}
\label{appendix:setups}

In this section, we provide additional details on experiments, including datasets, human evaluation setups, prompts (used for research idea generation and validation), and human-induced criteria. 


\subsection{Data Statistics}
We visualize a distribution of core paper categories used for idea generation in Figure~\ref{fig:discipline}, where the categories are obtained from Semantic Scholar API\footnote{https://www.semanticscholar.org/product/api}. From this, we find that the top 3 categories are computer science, medicine, and engineering. 

\subsection{Details on Human Evaluation}
To conduct evaluations with human judges, we recruited 10 researchers from the United States and South Korea, majoring in computer science, medicine, and biology, each with a minimum of 3 published papers. For annotation, they were provided with a 6-page guideline document, which includes the task instruction and annotation examples. After reading this document, the annotators access the Label Studio platform, on which they first read the title and abstract of the target paper, and then review and evaluate the generated research ideas from different models. During the evaluation process, they are allowed to use any external tools, such as web searches. We note that they were compensated at a rate of $\$22.20$ per hour. Also, on average, for one hour, they evaluated 3 sets of research ideas (that are generated from their own papers), with each set comprising three sub-ideas (the problem, method, and experiment design) from three different approaches (i.e., a total of 9 ideas for one hour). We perform three rounds of human evaluations with refinements in between, and, due to the cost associated with human annotations, we are able to fully evaluate a total of 150 ideas.

\begin{figure}[t!]
    \centering
    \includegraphics[width=0.975\columnwidth]{Figures/radar_model_norefine.pdf}
    \vspace{-0.1in}
    \caption{Results on our research idea generation task with model-based evaluation, where we exclude refinement steps.}
    \label{fig:norefine}
    \vspace{-0.1in}
\end{figure}

\subsection{Prompts for Ideas Generation}
\label{appendix:prompts:generation}

We provide the prompts used to elicit the idea generations from our full ResearchAgent, specifically for instantiating problem identification, method development, and experiment design in Table~\ref{tab:prompt:problem}, Table~\ref{tab:prompt:method}, and Table~\ref{tab:prompt:experiment}, respectively.

\subsection{Prompts for Idea Validation}
\label{appendix:prompts:validation}

We provide the prompts used to elicit the idea validation from our ReviewingAgents as well as the model-based evaluations, specifically for instantiating problem validation, method validation, and experiment design validation in Table~\ref{tab:prompt:problem:valid}, Table~\ref{tab:prompt:method:valid}, and Table~\ref{tab:prompt:experiment:valid}, respectively. In addition, we provide the criteria used, which are induced by human judgments in the next subsection (Appendix~\ref{appendix:criteria}).

\subsection{Criteria Induced by Human Judgements}
\label{appendix:criteria}

Recall that, to align model-based evaluations with human preferences, we induce the criteria (used for automatic evaluations) with actual human judgments. We note that this is done by prompting GPT-4 with 10 pairs of generated ideas and (randomly selected) human judgments. We provide the resulting criteria for validations of problems, methods, and experiment designs in Table~\ref{tab:criteria:problem}, Table~\ref{tab:criteria:method}, and Table~\ref{tab:criteria:experiment}, respectively.  



\section{Additional Experimental Results}
We provide additional experimental results, including comparisons without refinements and examples of the generated research ideas. 


\subsection{Results without Refinement Steps}
To see whether the proposed ResearchAgent is consistently effective even without ReviewingAgents, we show the model-based evaluation results without any refinement steps in Figure~\ref{fig:norefine}. From this, we clearly observe that the full ResearchAgent outperforms its variants, demonstrating its effectiveness.

\begin{figure*}[t!]
    \centering
    \includegraphics[width=0.9\linewidth]{Figures/radar_model_domain.pdf}
    \vspace{-0.1in}
    \caption{Breakdown results of the research ideas generated from our full ResearchAgent across different domains.}
    \label{fig:domain}
    \vspace{-0.075in}
\end{figure*}

\subsection{Results on Generated Ideas by Domain}
To see the quality of the generated research ideas across different domains, we breakdown the performance of ResearchAgent according to the categories of core papers in Figure~\ref{fig:discipline}, and present the results in Figure~\ref{fig:domain}. From this, we observe that the generated research ideas on the high-resource domains (such as Computer Science, Medicine, and Engineering where there is a greater volume of existing literature as shown in Figure~\ref{fig:discipline}) are superior to those generated from the low-resource domain papers (such as Physics, Chemistry, and Mathematics). This disparity might be attributed to the fact that the underlying LLMs used to generate research ideas are likely trained on data predominantly sourced from high-resource domains, which leads to enhancing their ability to comprehend scientific concepts and produce relevant research ideas.

\begin{table}[t!]
\caption{Results with different entity retrieval strategies.}
\vspace{-0.1in}
\label{tab:retrieval}
\small
\centering
\resizebox{0.475\textwidth}{!}{
\renewcommand{\arraystretch}{1.15}
% \renewcommand{\tabcolsep}{2.0mm}
\begin{tabular}{lccc}
\toprule

\textbf{Methods} & \textbf{Problem} & \textbf{Method} & \textbf{Experiment} \\

\midrule
\midrule

ResearchAgent &  \\

\noalign{\vskip 0.25ex}\cdashline{1-4}\noalign{\vskip 0.75ex}

- w/ Co-occurrence-based Retrieval & \textbf{4.52 }& 4.28 & \textbf{4.18} \\

- w/ Embedding-based Retrieval & 4.49 & \textbf{4.34} & 4.16 \\

- w/o Entity Retrieval & 4.35 & 4.13 & 4.02 \\

\bottomrule

\end{tabular}
}
\vspace{-0.05in}
\end{table}


\subsection{Analysis with Different Entity Retrieval}
\label{appendix:retrieval}

To see the effectiveness of different entity retrieval strategies, we perform additional experiments, replacing the co-occurrence-based entity retrieval in Equation~\ref{eq:retrieval:detail} to the contextual embedding-based retrieval. Notably, this contextual embedding-based retrieval approach uses the entities that have the highest similarity to the entities appearing in the current literature (i.e., core paper and its references) used for idea generation, where the similarities are calculated based on embedding-level similarities between entities over the latent space represented by the entity linker~\cite{blink}. Therefore, unlike the previous co-occurrence-based entity retrieval that may retrieve entities that have opposite concepts to the main idea of the current core paper (since we often mention limitations of previous work along with the proposed ideas), this embedding-based approach may provide the ResearchAgent with mostly the entities having similar concepts to the core paper. Nevertheless, as shown in Table~\ref{tab:retrieval}, the results with the strategy of entity co-occurrence-based retrieval are comparable to the results with the new embedding-based contextual retrieval. These results might confirm that there is not much difference in the quality of entity retrieval among those two strategies, i.e., most entities retrieved from the co-occurrence-based retrieval are contextually relevant for generating research ideas. 


\begin{table*}
    \caption{The prompt used in the full instantiation of ResearchAgent for problem identification.}
    \label{tab:prompt:problem}
    \vspace{-0.1in}
    \resizebox{0.95\textwidth}{!}{
    \renewcommand{\arraystretch}{1.1}
    \renewcommand{\tabcolsep}{2.5mm}
        \begin{tabular}{ll}
        \toprule
        \multicolumn{1}{p{.18\textwidth}}{\textbf{Types}} & \makecell{\multicolumn{1}{p{.92\textwidth}}{\textbf{Texts}}} \\
        \midrule
        \multicolumn{1}{p{.18\textwidth}}{\textbf{System Message}} & 
        \makecell{
            \multicolumn{1}{p{.92\textwidth}}{You are an AI assistant whose primary goal is to identify promising, new, and key scientific problems based on existing scientific literature, in order to aid researchers in discovering novel and significant research opportunities that can advance the field.}
        } \\
        \midrule
        \multicolumn{1}{p{.18\textwidth}}{\textbf{User Message}} & 
        \makecell{
            \multicolumn{1}{p{.92\textwidth}}{{You are going to generate a research problem that should be original, clear, feasible, relevant, and significant to its field. This will be based on the title and abstract of the target paper, those of \{len(references)\} related papers in the existing literature, and \{len(entities)\} entities potentially connected to the research area.}} \\ \\
            \multicolumn{1}{p{.92\textwidth}}{{Understanding of the target paper, related papers, and entities is essential:}} \\
            \multicolumn{1}{p{.92\textwidth}}{{- The target paper is the primary research study you aim to enhance or build upon through future research, serving as the central source and focus for identifying and developing the specific research problem.}} \\
            \multicolumn{1}{p{.92\textwidth}}{{- The related papers are studies that have cited the target paper, indicating their direct relevance and connection to the primary research topic you are focusing on, and providing additional context and insights that are essential for understanding and expanding upon the target paper.}} \\
            \multicolumn{1}{p{.92\textwidth}}{{- The entities can include topics, keywords, individuals, events, or any subjects with possible direct or indirect connections to the target paper or the related studies, serving as auxiliary sources of inspiration or information that may be instrumental in formulating the research problem.}} \\ \\
            \multicolumn{1}{p{.92\textwidth}}{{Your approach should be systematic:}} \\ 
            \multicolumn{1}{p{.92\textwidth}}{{- Start by thoroughly reading the title and abstract of the target paper to understand its core focus.}} \\
            \multicolumn{1}{p{.92\textwidth}}{{- Next, proceed to read the titles and abstracts of the related papers to gain a broader perspective and insights relevant to the primary research topic.}} \\
            \multicolumn{1}{p{.92\textwidth}}{{- Finally, explore the entities to further broaden your perspective, drawing upon a diverse pool of inspiration and information, while keeping in mind that not all may be relevant.}} \\ \\
            \multicolumn{1}{p{.92\textwidth}}{{I am going to provide the target paper, related papers, and entities, as follows: }} \\
            \multicolumn{1}{p{.92\textwidth}}{{Target paper title: \{paper['title']\}}} \\
            \multicolumn{1}{p{.92\textwidth}}{{Target paper abstract: \{paper['abstract']\} }} \\
            \multicolumn{1}{p{.92\textwidth}}{{Related paper titles: \{relatedPaper['titles']\}}} \\
            \multicolumn{1}{p{.92\textwidth}}{{Related paper abstracts: \{relatedPaper['abstracts']\} }} \\
            \multicolumn{1}{p{.92\textwidth}}{{Entities: \{Entities\} }} \\ \\
            \multicolumn{1}{p{.92\textwidth}}{{With the provided target paper, related papers, and entities, your objective now is to formulate a research problem that not only builds upon these existing studies but also strives to be original, clear, feasible, relevant, and significant. Before crafting the research problem, revisit the title and abstract of the target paper, to ensure it remains the focal point of your research problem identification process.}} \\ \\
            \multicolumn{1}{p{.92\textwidth}}{{Target paper title: \{paper['title']\}}} \\
            \multicolumn{1}{p{.92\textwidth}}{{Target paper abstract: \{paper['abstract']\} }} \\ \\
            \multicolumn{1}{p{.92\textwidth}}{{Then, following your review of the above content, please proceed to generate one research problem with the rationale, in the format of }} \\
            \multicolumn{1}{p{.92\textwidth}}{{Problem: }} \\
            \multicolumn{1}{p{.92\textwidth}}{{Rationale: }} \\
        } \\
        \bottomrule
        \end{tabular}
    }
\end{table*}

\begin{table*}
    \caption{The prompt used in the full instantiation of ResearchAgent for method development.}
    \label{tab:prompt:method}
    \vspace{-0.1in}
    \resizebox{0.95\textwidth}{!}{
    \renewcommand{\arraystretch}{1.1}
    \renewcommand{\tabcolsep}{2.5mm}
        \begin{tabular}{ll}
        \toprule
        \multicolumn{1}{p{.18\textwidth}}{\textbf{Types}} & \makecell{\multicolumn{1}{p{.92\textwidth}}{\textbf{Texts}}} \\
        \midrule
        \multicolumn{1}{p{.18\textwidth}}{\textbf{System Message}} & 
        \makecell{
            \multicolumn{1}{p{.92\textwidth}}{You are an AI assistant whose primary goal is to propose innovative, rigorous, and valid methodologies to solve newly identified scientific problems derived from existing scientific literature, in order to empower researchers to pioneer groundbreaking solutions that catalyze breakthroughs in their fields.}
        } \\
        \midrule
        \multicolumn{1}{p{.18\textwidth}}{\textbf{User Message}} & 
        \makecell{
            \multicolumn{1}{p{.92\textwidth}}{{You are going to propose a scientific method to address a specific research problem. Your method should be clear, innovative, rigorous, valid, and generalizable. This will be based on a deep understanding of the research problem, its rationale, existing studies, and various entities.}} \\ \\
            \multicolumn{1}{p{.92\textwidth}}{{Understanding of the research problem, existing studies, and entities is essential:}} \\
            \multicolumn{1}{p{.92\textwidth}}{{- The research problem has been formulated based on an in-depth review of existing studies and a potential exploration of relevant entities, which should be the cornerstone of your method development.}} \\
            \multicolumn{1}{p{.92\textwidth}}{{- The existing studies refer to the target paper that has been pivotal in identifying the problem, as well as the related papers that have been additionally referenced in the problem discovery phase, all serving as foundational material for developing the method.}} \\
            \multicolumn{1}{p{.92\textwidth}}{{- The entities can include topics, keywords, individuals, events, or any subjects with possible direct or indirect connections to the existing studies, serving as auxiliary sources of inspiration or information that may be instrumental in method development.}} \\ \\
            \multicolumn{1}{p{.92\textwidth}}{{Your approach should be systematic:}} \\ 
            \multicolumn{1}{p{.92\textwidth}}{{- Start by thoroughly reading the research problem and its rationale, to understand your primary focus.}} \\
            \multicolumn{1}{p{.92\textwidth}}{{- Next, proceed to review the titles and abstracts of existing studies, to gain a broader perspective and insights relevant to the primary research topic.}} \\
            \multicolumn{1}{p{.92\textwidth}}{{- Finally, explore the entities to further broaden your perspective, drawing upon a diverse pool of inspiration and information, while keeping in mind that not all may be relevant.}} \\ \\
            \multicolumn{1}{p{.92\textwidth}}{{I am going to provide the research problem, existing studies (target paper \& related papers), and entities, as follows: }} \\
            \multicolumn{1}{p{.92\textwidth}}{{Research problem: \{researchProblem\}}} \\
            \multicolumn{1}{p{.92\textwidth}}{{Rationale: \{researchProblemRationale\}}} \\
            \multicolumn{1}{p{.92\textwidth}}{{Target paper title: \{paper['title']\}}} \\
            \multicolumn{1}{p{.92\textwidth}}{{Target paper abstract: \{paper['abstract']\} }} \\
            \multicolumn{1}{p{.92\textwidth}}{{Related paper titles: \{relatedPaper['titles']\}}} \\
            \multicolumn{1}{p{.92\textwidth}}{{Related paper abstracts: \{relatedPaper['abstracts']\} }} \\
            \multicolumn{1}{p{.92\textwidth}}{{Entities: \{Entities\} }} \\ \\
            \multicolumn{1}{p{.92\textwidth}}{{With the provided research problem, existing studies, and entities, your objective now is to formulate a method that not only leverages these resources but also strives to be clear, innovative, rigorous, valid, and generalizable. Before crafting the method, revisit the research problem, to ensure it remains the focal point of your method development process.}} \\ \\
            \multicolumn{1}{p{.92\textwidth}}{{Research problem: \{researchProblem\}}} \\
            \multicolumn{1}{p{.92\textwidth}}{{Rationale: \{researchProblemRationale\}}} \\ \\
            \multicolumn{1}{p{.92\textwidth}}{{Then, following your review of the above content, please proceed to propose your method with its rationale, in the format of }} \\
            \multicolumn{1}{p{.92\textwidth}}{{Method: }} \\
            \multicolumn{1}{p{.92\textwidth}}{{Rationale: }} \\
        } \\
        \bottomrule
        \end{tabular}
    }
\end{table*}

\begin{table*}
    \caption{The prompt used in the full instantiation of ResearchAgent for experiment design.}
    \label{tab:prompt:experiment}
    \vspace{-0.1in}
    \resizebox{0.95\textwidth}{!}{
    \renewcommand{\arraystretch}{1.1}
    \renewcommand{\tabcolsep}{2.5mm}
        \begin{tabular}{ll}
        \toprule
        \multicolumn{1}{p{.18\textwidth}}{\textbf{Types}} & \makecell{\multicolumn{1}{p{.92\textwidth}}{\textbf{Texts}}} \\
        \midrule
        \multicolumn{1}{p{.18\textwidth}}{\textbf{System Message}} & 
        \makecell{
            \multicolumn{1}{p{.92\textwidth}}{You are an AI assistant whose primary goal is to design robust, feasible, and impactful experiments based on identified scientific problems and proposed methodologies from existing scientific literature, in order to enable researchers to systematically test hypotheses and validate groundbreaking discoveries that can transform their respective fields.}
        } \\
        \midrule
        \multicolumn{1}{p{.18\textwidth}}{\textbf{User Message}} & 
        \makecell{
            \multicolumn{1}{p{.92\textwidth}}{{You are going to design an experiment, aimed at validating a proposed method to address a specific research problem. Your experiment design should be clear, robust, reproducible, valid, and feasible. This will be based on a deep understanding of the research problem, scientific method, existing studies, and various entities.}} \\ \\
            \multicolumn{1}{p{.92\textwidth}}{{Understanding of the research problem, scientific method, existing studies, and entities is essential:}} \\
            \multicolumn{1}{p{.92\textwidth}}{{- The research problem has been formulated based on an in-depth review of existing studies and a potential exploration of relevant entities.}} \\
            \multicolumn{1}{p{.92\textwidth}}{{- The scientific method has been proposed to tackle the research problem, which has been informed by insights gained from existing studies and relevant entities.}} \\
            \multicolumn{1}{p{.92\textwidth}}{{- The existing studies refer to the target paper that has been pivotal in identifying the problem and method, as well as the related papers that have been additionally referenced in the discovery phase of the problem and method, all serving as foundational material for designing the experiment.}} \\
            \multicolumn{1}{p{.92\textwidth}}{{- The entities can include topics, keywords, individuals, events, or any subjects with possible direct or indirect connections to the existing studies, serving as auxiliary sources of inspiration or information that may be instrumental in your experiment design.}} \\ \\
            \multicolumn{1}{p{.92\textwidth}}{{Your approach should be systematic:}} \\ 
            \multicolumn{1}{p{.92\textwidth}}{{- Start by thoroughly reading the research problem and its rationale followed by the proposed method and its rationale, to pinpoint your primary focus.}} \\
            \multicolumn{1}{p{.92\textwidth}}{{- Next, proceed to review the titles and abstracts of existing studies, to gain a broader perspective and insights relevant to the primary research topic.}} \\
            \multicolumn{1}{p{.92\textwidth}}{{- Finally, explore the entities to further broaden your perspective, drawing upon a diverse pool of inspiration and information, while keeping in mind that not all may be relevant.}} \\ \\
            \multicolumn{1}{p{.92\textwidth}}{{I am going to provide the research problem, scientific method, existing studies (target paper \& related papers), and entities, as follows: }} \\
            \multicolumn{1}{p{.92\textwidth}}{{Research problem: \{researchProblem\}}} \\
            \multicolumn{1}{p{.92\textwidth}}{{Rationale: \{researchProblemRationale\}}} \\
            \multicolumn{1}{p{.92\textwidth}}{{Scientific method: \{scientificMethod\}}} \\
            \multicolumn{1}{p{.92\textwidth}}{{Rationale: \{scientificMethodRationale\}}} \\
            \multicolumn{1}{p{.92\textwidth}}{{Target paper title: \{paper['title']\}}} \\
            \multicolumn{1}{p{.92\textwidth}}{{Target paper abstract: \{paper['abstract']\} }} \\
            \multicolumn{1}{p{.92\textwidth}}{{Related paper titles: \{relatedPaper['titles']\}}} \\
            \multicolumn{1}{p{.92\textwidth}}{{Related paper abstracts: \{relatedPaper['abstracts']\} }} \\
            \multicolumn{1}{p{.92\textwidth}}{{Entities: \{Entities\} }} \\ \\
            \multicolumn{1}{p{.92\textwidth}}{{With the provided research problem, scientific method, existing studies, and entities, your objective now is to design an experiment that not only leverages these resources but also strives to be clear, robust, reproducible, valid, and feasible. Before crafting the experiment design, revisit the research problem and proposed method, to ensure they remain at the center of your experiment design process.}} \\ \\
            \multicolumn{1}{p{.92\textwidth}}{{Research problem: \{researchProblem\}}} \\
            \multicolumn{1}{p{.92\textwidth}}{{Rationale: \{researchProblemRationale\}}} \\ 
            \multicolumn{1}{p{.92\textwidth}}{{Scientific method: \{scientificMethod\}}} \\
            \multicolumn{1}{p{.92\textwidth}}{{Rationale: \{scientificMethodRationale\}}} \\ \\
            \multicolumn{1}{p{.92\textwidth}}{{Then, following your review of the above content, please proceed to outline your experiment with its rationale, in the format of }} \\
            \multicolumn{1}{p{.92\textwidth}}{{Experiment: }} \\
            \multicolumn{1}{p{.92\textwidth}}{{Rationale: }} \\
        } \\
        \bottomrule
        \end{tabular}
    }
\end{table*}
\begin{table*}
    \caption{The prompt used in the full instantiation of ReviewingAgent for problem validation.}
    \label{tab:prompt:problem:valid}
    \vspace{-0.1in}
    \resizebox{0.95\textwidth}{!}{
    \renewcommand{\arraystretch}{1.1}
    \renewcommand{\tabcolsep}{2.5mm}
        \begin{tabular}{ll}
        \toprule
        \multicolumn{1}{p{.18\textwidth}}{\textbf{Types}} & \makecell{\multicolumn{1}{p{.92\textwidth}}{\textbf{Texts}}} \\
        \midrule
        \multicolumn{1}{p{.18\textwidth}}{\textbf{System Message}} & 
        \makecell{
            \multicolumn{1}{p{.92\textwidth}}{You are an AI assistant whose primary goal is to assess the quality and validity of scientific problems across diverse dimensions, in order to aid researchers in refining their problems based on your evaluations and feedback, thereby enhancing the impact and reach of their work.}
        } \\
        \midrule
        \multicolumn{1}{p{.18\textwidth}}{\textbf{User Message}} & 
        \makecell{
            \multicolumn{1}{p{.92\textwidth}}{{You are going to evaluate a research problem for its \{metric\}, focusing on how well it is defined in a clear, precise, and understandable manner.}} \\ \\
            \multicolumn{1}{p{.92\textwidth}}{{As part of your evaluation, you can refer to the existing studies that may be related to the problem, which will help in understanding the context of the problem for a more comprehensive assessment.}} \\
            \multicolumn{1}{p{.92\textwidth}}{{- The existing studies refer to the target paper that has been pivotal in identifying the problem, as well as the related papers that have been additionally referenced in the discovery phase of the problem.}} \\ \\
            \multicolumn{1}{p{.92\textwidth}}{{The existing studies (target paper \& related papers) are as follows:}} \\ 
            \multicolumn{1}{p{.92\textwidth}}{{Target paper title: \{paper['title']\}}} \\
            \multicolumn{1}{p{.92\textwidth}}{{Target paper abstract: \{paper['abstract']\} }} \\
            \multicolumn{1}{p{.92\textwidth}}{{Related paper titles: \{relatedPaper['titles']\}}} \\
            \multicolumn{1}{p{.92\textwidth}}{{Related paper abstracts: \{relatedPaper['abstracts']\} }} \\ \\
            \multicolumn{1}{p{.92\textwidth}}{{Now, proceed with your \{metric\} evaluation approach that should be systematic:}} \\
            \multicolumn{1}{p{.92\textwidth}}{{- Start by thoroughly reading the research problem and its rationale, keeping in mind the context provided by the existing studies mentioned above.}} \\
            \multicolumn{1}{p{.92\textwidth}}{{- Next, generate a review and feedback that should be constructive, helpful, and concise, focusing on the \{metric\} of the problem.}} \\
            \multicolumn{1}{p{.92\textwidth}}{{- Finally, provide a score on a 5-point Likert scale, with 1 being the lowest, please ensuring a discerning and critical evaluation to avoid a tendency towards uniformly high ratings (4-5) unless fully justified:}} \\ 
            \multicolumn{1}{p{.92\textwidth}}{{\{criteria\}}} \\ \\
            \multicolumn{1}{p{.92\textwidth}}{{I am going to provide the research problem with its rationale, as follows:}} \\
            \multicolumn{1}{p{.92\textwidth}}{{Research problem: \{researchProblem\}}} \\
            \multicolumn{1}{p{.92\textwidth}}{{Rationale: \{researchProblemRationale\}}} \\ \\
            \multicolumn{1}{p{.92\textwidth}}{{After your evaluation of the above content, please provide your review, feedback, and rating, in the format of}} \\ 
            \multicolumn{1}{p{.92\textwidth}}{{Review: }} \\
            \multicolumn{1}{p{.92\textwidth}}{{Feedback: }} \\
            \multicolumn{1}{p{.92\textwidth}}{{Rating (1-5): }} \\
        } \\
        \bottomrule
        \end{tabular}
    }
\end{table*}

\begin{table*}
    \caption{The prompt used in the full instantiation of ReviewingAgent for method validation.}
    \label{tab:prompt:method:valid}
    \vspace{-0.1in}
    \resizebox{0.95\textwidth}{!}{
    \renewcommand{\arraystretch}{1.1}
    \renewcommand{\tabcolsep}{2.5mm}
        \begin{tabular}{ll}
        \toprule
        \multicolumn{1}{p{.18\textwidth}}{\textbf{Types}} & \makecell{\multicolumn{1}{p{.92\textwidth}}{\textbf{Texts}}} \\
        \midrule
        \multicolumn{1}{p{.18\textwidth}}{\textbf{System Message}} & 
        \makecell{
            \multicolumn{1}{p{.92\textwidth}}{You are an AI assistant whose primary goal is to assess the quality and soundness of scientific methods across diverse dimensions, in order to aid researchers in refining their methods based on your evaluations and feedback, thereby enhancing the impact and reach of their work.}
        } \\
        \midrule
        \multicolumn{1}{p{.18\textwidth}}{\textbf{User Message}} & 
        \makecell{
            \multicolumn{1}{p{.92\textwidth}}{{You are going to evaluate a scientific method for its \{metric\} in addressing a research problem, focusing on how well it is described in a clear, precise, and understandable manner that allows for replication and comprehension of the approach.}} \\ \\
            \multicolumn{1}{p{.92\textwidth}}{{As part of your evaluation, you can refer to the research problem, and existing studies, which will help in understanding the context of the proposed method for a more comprehensive assessment.}} \\
            \multicolumn{1}{p{.92\textwidth}}{{- The research problem has been used as the cornerstone of the method development, formulated based on an in-depth review of existing studies and a potential exploration of relevant entities.}} \\
            \multicolumn{1}{p{.92\textwidth}}{{- The existing studies refer to the target paper that has been pivotal in identifying the problem and method, as well as the related papers that have been additionally referenced in the discovery phase of the problem and method.}} \\ \\
            \multicolumn{1}{p{.92\textwidth}}{{The research problem and existing studies (target paper \& related papers) are as follows:}} \\ 
            \multicolumn{1}{p{.92\textwidth}}{{Research problem: \{researchProblem\}}} \\
            \multicolumn{1}{p{.92\textwidth}}{{Rationale: \{researchProblemRationale\}}} \\ 
            \multicolumn{1}{p{.92\textwidth}}{{Target paper title: \{paper['title']\}}} \\
            \multicolumn{1}{p{.92\textwidth}}{{Target paper abstract: \{paper['abstract']\} }} \\
            \multicolumn{1}{p{.92\textwidth}}{{Related paper titles: \{relatedPaper['titles']\}}} \\
            \multicolumn{1}{p{.92\textwidth}}{{Related paper abstracts: \{relatedPaper['abstracts']\} }} \\ \\
            \multicolumn{1}{p{.92\textwidth}}{{Now, proceed with your \{metric\} evaluation approach that should be systematic:}} \\
            \multicolumn{1}{p{.92\textwidth}}{{- Start by thoroughly reading the proposed method and its rationale, keeping in mind the context provided by the research problem, and existing studies mentioned above.}} \\
            \multicolumn{1}{p{.92\textwidth}}{{- Next, generate a review and feedback that should be constructive, helpful, and concise, focusing on the \{metric\} of the method.}} \\
            \multicolumn{1}{p{.92\textwidth}}{{- Finally, provide a score on a 5-point Likert scale, with 1 being the lowest, please ensuring a discerning and critical evaluation to avoid a tendency towards uniformly high ratings (4-5) unless fully justified:}} \\ 
            \multicolumn{1}{p{.92\textwidth}}{{\{criteria\}}} \\ \\
            \multicolumn{1}{p{.92\textwidth}}{{I am going to provide the proposed method with its rationale, as follows:}} \\
            \multicolumn{1}{p{.92\textwidth}}{{Scientific method: \{scientificMethod\}}} \\
            \multicolumn{1}{p{.92\textwidth}}{{Rationale: \{scientificMethodRationale\}}} \\ \\
            \multicolumn{1}{p{.92\textwidth}}{{After your evaluation of the above content, please provide your review, feedback, and rating, in the format of}} \\ 
            \multicolumn{1}{p{.92\textwidth}}{{Review: }} \\
            \multicolumn{1}{p{.92\textwidth}}{{Feedback: }} \\
            \multicolumn{1}{p{.92\textwidth}}{{Rating (1-5): }} \\
        } \\
        \bottomrule
        \end{tabular}
    }
\end{table*}

\begin{table*}
    \caption{The prompt used in the full instantiation of ReviewingAgent for experiment design validation.}
    \label{tab:prompt:experiment:valid}
    \vspace{-0.1in}
    \resizebox{0.95\textwidth}{!}{
    \renewcommand{\arraystretch}{1.1}
    \renewcommand{\tabcolsep}{2.5mm}
        \begin{tabular}{ll}
        \toprule
        \multicolumn{1}{p{.18\textwidth}}{\textbf{Types}} & \makecell{\multicolumn{1}{p{.92\textwidth}}{\textbf{Texts}}} \\
        \midrule
        \multicolumn{1}{p{.18\textwidth}}{\textbf{System Message}} & 
        \makecell{
            \multicolumn{1}{p{.92\textwidth}}{You are an AI assistant whose primary goal is to meticulously evaluate the experimental designs of scientific papers across diverse dimensions, in order to aid researchers in refining their experimental approaches based on your evaluations and feedback, thereby amplifying the quality and impact of their scientific contributions.}
        } \\
        \midrule
        \multicolumn{1}{p{.18\textwidth}}{\textbf{User Message}} & 
        \makecell{
            \multicolumn{1}{p{.92\textwidth}}{{You are going to evaluate an experiment design for its \{metric\} in validating a scientific method to address a research problem, focusing on how well it is described in a clear, precise, and understandable manner, enabling others to grasp the setup, procedure, and expected outcomes.}} \\ \\
            \multicolumn{1}{p{.92\textwidth}}{{As part of your evaluation, you can refer to the research problem, scientific method, and existing studies, which will help in understanding the context of the designed experiment for a more comprehensive assessment.}} \\
            \multicolumn{1}{p{.92\textwidth}}{{- The research problem has been formulated based on an in-depth review of existing studies and a potential exploration of relevant entities.}} \\
            \multicolumn{1}{p{.92\textwidth}}{{- The scientific method has been proposed to tackle the research problem, which has been informed by insights gained from existing studies and relevant entities.}} \\ 
            \multicolumn{1}{p{.92\textwidth}}{{- The existing studies refer to the target paper that has been pivotal in identifying the problem, method, and experiment, as well as the related papers that have been additionally referenced in their discovery phases.}} \\ \\
            \multicolumn{1}{p{.92\textwidth}}{{The research problem, scientific method, and existing studies (target paper \& related papers) are as follows:}} \\ 
            \multicolumn{1}{p{.92\textwidth}}{{Research problem: \{researchProblem\}}} \\
            \multicolumn{1}{p{.92\textwidth}}{{Rationale: \{researchProblemRationale\}}} \\ 
            \multicolumn{1}{p{.92\textwidth}}{{Scientific method: \{scientificMethod\}}} \\
            \multicolumn{1}{p{.92\textwidth}}{{Rationale: \{scientificMethodRationale\}}} \\
            \multicolumn{1}{p{.92\textwidth}}{{Target paper title: \{paper['title']\}}} \\
            \multicolumn{1}{p{.92\textwidth}}{{Target paper abstract: \{paper['abstract']\} }} \\
            \multicolumn{1}{p{.92\textwidth}}{{Related paper titles: \{relatedPaper['titles']\}}} \\
            \multicolumn{1}{p{.92\textwidth}}{{Related paper abstracts: \{relatedPaper['abstracts']\} }} \\ \\
            \multicolumn{1}{p{.92\textwidth}}{{Now, proceed with your \{metric\} evaluation approach that should be systematic:}} \\
            \multicolumn{1}{p{.92\textwidth}}{{- Start by thoroughly reading the experiment design and its rationale, keeping in mind the context provided by the research problem, scientific method, and existing studies mentioned above.}} \\
            \multicolumn{1}{p{.92\textwidth}}{{- Next, generate a review and feedback that should be constructive, helpful, and concise, focusing on the \{metric\} of the experiment.}} \\
            \multicolumn{1}{p{.92\textwidth}}{{- Finally, provide a score on a 5-point Likert scale, with 1 being the lowest, please ensuring a discerning and critical evaluation to avoid a tendency towards uniformly high ratings (4-5) unless fully justified:}} \\ 
            \multicolumn{1}{p{.92\textwidth}}{{\{criteria\}}} \\ \\
            \multicolumn{1}{p{.92\textwidth}}{{I am going to provide the designed experiment with its rationale, as follows:}} \\
            \multicolumn{1}{p{.92\textwidth}}{{Experiment design: \{experimentDesign\}}} \\
            \multicolumn{1}{p{.92\textwidth}}{{Rationale: \{experimentDesignRationale\}}} \\ \\
            \multicolumn{1}{p{.92\textwidth}}{{After your evaluation of the above content, please provide your review, feedback, and rating, in the format of}} \\ 
            \multicolumn{1}{p{.92\textwidth}}{{Review: }} \\
            \multicolumn{1}{p{.92\textwidth}}{{Feedback: }} \\
            \multicolumn{1}{p{.92\textwidth}}{{Rating (1-5): }} \\
        } \\
        \bottomrule
        \end{tabular}
    }
\end{table*}

\begin{table*}
    \caption{The criteria used for evaluating research ideas: problems, methods, and experiment designs.}
    \label{tab:criteria}
    \vspace{-0.1in}
    \resizebox{0.95\textwidth}{!}{
    \renewcommand{\arraystretch}{1.1}
    \renewcommand{\tabcolsep}{2.5mm}
        \begin{tabular}{lll}
        \toprule
        \multicolumn{1}{p{.12\textwidth}}{\textbf{Types}} & \makecell{\multicolumn{1}{p{.12\textwidth}}{\textbf{Criteria}}} & \makecell{\multicolumn{1}{p{1.05\textwidth}}{\textbf{Texts}}} \\
        \midrule
        \multirow{8}{*}{\textbf{Problem}} 
        & \makecell{\multicolumn{1}{p{.12\textwidth}}{\textbf{Clarity}}} & \makecell{
            \multicolumn{1}{p{1.05\textwidth}}{It assesses whether the problem is defined in a clear, precise, and understandable manner.}
        } \\
        \noalign{\vskip 0.25ex}\cdashline{2-3}\noalign{\vskip 0.75ex}
        & \makecell{\multicolumn{1}{p{.12\textwidth}}{\textbf{Relevance}}} & \makecell{
            \multicolumn{1}{p{1.05\textwidth}}{It measures whether the problem is pertinent and applicable to the current field or context of study.}
        } \\
        \noalign{\vskip 0.25ex}\cdashline{2-3}\noalign{\vskip 0.75ex}
        & \makecell{\multicolumn{1}{p{.12\textwidth}}{\textbf{Originality}}} & \makecell{
            \multicolumn{1}{p{1.05\textwidth}}{It evaluates whether the problem presents a novel challenge or unique perspective that has not been extensively explored before.}
        } \\
        \noalign{\vskip 0.25ex}\cdashline{2-3}\noalign{\vskip 0.75ex}
        & \makecell{\multicolumn{1}{p{.12\textwidth}}{\textbf{Feasibility}}} & \makecell{
            \multicolumn{1}{p{1.05\textwidth}}{It examines whether the problem can realistically be investigated or solved with the available resources and within reasonable constraints.}
        } \\
        \noalign{\vskip 0.25ex}\cdashline{2-3}\noalign{\vskip 0.75ex}
        & \makecell{\multicolumn{1}{p{.12\textwidth}}{\textbf{Significance}}} & \makecell{
            \multicolumn{1}{p{1.05\textwidth}}{It assesses the importance and potential impact of solving the problem, including its contribution to the field or its broader implications.}
        } \\
        \midrule
        \multirow{10}{*}{\textbf{Method}} 
        & \makecell{\multicolumn{1}{p{.12\textwidth}}{\textbf{Clarity}}} & \makecell{
            \multicolumn{1}{p{1.05\textwidth}}{It assesses whether the method is described in a clear, precise, and understandable manner that allows for replication and comprehension of the approach.}
        } \\
        \noalign{\vskip 0.25ex}\cdashline{2-3}\noalign{\vskip 0.75ex}
        & \makecell{\multicolumn{1}{p{.12\textwidth}}{\textbf{Validity}}} & \makecell{
            \multicolumn{1}{p{1.05\textwidth}}{It measures the accuracy, relevance, and soundness of the method in addressing the research problem, ensuring that it is appropriate and directly relevant to the objectives of the study.}
        } \\
        \noalign{\vskip 0.25ex}\cdashline{2-3}\noalign{\vskip 0.75ex}
        & \makecell{\multicolumn{1}{p{.12\textwidth}}{\textbf{Rigorousness}}} & \makecell{
            \multicolumn{1}{p{1.05\textwidth}}{It examines the thoroughness, precision, and consistency of the method, ensuring that the approach is systematic, well-structured, and adheres to high standards of research quality.}
        } \\
        \noalign{\vskip 0.25ex}\cdashline{2-3}\noalign{\vskip 0.75ex}
        & \makecell{\multicolumn{1}{p{.12\textwidth}}{\textbf{Innovativeness}}} & \makecell{
            \multicolumn{1}{p{1.05\textwidth}}{It evaluates whether the method introduces new techniques, approaches, or perspectives to the research field that differ from standard research practices and advance them in the field.}
        } \\
        \noalign{\vskip 0.25ex}\cdashline{2-3}\noalign{\vskip 0.75ex}
        & \makecell{\multicolumn{1}{p{.12\textwidth}}{\textbf{Generalizability}}} & \makecell{
            \multicolumn{1}{p{1.05\textwidth}}{It assesses the extent to which the method can be applied to or is relevant for other contexts, populations, or settings beyond the scope of the study.}
        } \\
        \midrule
        \multirow{12}{*}{\textbf{Experiment}} 
        & \makecell{\multicolumn{1}{p{.12\textwidth}}{\textbf{Clarity}}} & \makecell{
            \multicolumn{1}{p{1.05\textwidth}}{It determines whether the experiment design is described in a clear, precise, and understandable manner, enabling others to grasp the setup, procedure, and expected outcomes.}
        } \\
        \noalign{\vskip 0.25ex}\cdashline{2-3}\noalign{\vskip 0.75ex}
        & \makecell{\multicolumn{1}{p{.12\textwidth}}{\textbf{Validity}}} & \makecell{
            \multicolumn{1}{p{1.05\textwidth}}{It measures the appropriateness and soundness of the experimental design in accurately addressing the research questions or effectively validating the proposed methods, ensuring that the design effectively tests what it is intended to examine.}
        } \\
        \noalign{\vskip 0.25ex}\cdashline{2-3}\noalign{\vskip 0.75ex}
        & \makecell{\multicolumn{1}{p{.12\textwidth}}{\textbf{Robustness}}} & \makecell{
            \multicolumn{1}{p{1.05\textwidth}}{It evaluates the durability of the experimental design across a wide range of conditions and variables, ensuring that the outcomes are not reliant on a few specific cases and remain consistent across a broad spectrum of scenarios.}
        } \\
        \noalign{\vskip 0.25ex}\cdashline{2-3}\noalign{\vskip 0.75ex}
        & \makecell{\multicolumn{1}{p{.12\textwidth}}{\textbf{Feasibility}}} & \makecell{
            \multicolumn{1}{p{1.05\textwidth}}{It evaluates whether the experiment design can realistically be implemented with the available resources, time, and technological or methodological constraints, ensuring that the experiment is practical and achievable.}
        } \\
        \noalign{\vskip 0.25ex}\cdashline{2-3}\noalign{\vskip 0.75ex}
        & \makecell{\multicolumn{1}{p{.12\textwidth}}{\textbf{Reproducibility}}} & \makecell{
            \multicolumn{1}{p{1.05\textwidth}}{It examines whether the information provided is sufficient and detailed enough for other researchers to reproduce the experiment using the same methodology and conditions, ensuring the reliability of the findings.}
        } \\
        \bottomrule
        \end{tabular}
    }
\end{table*}

\begin{table*}
    \caption{The criteria induced from human judgments for validating the identified problems, which are used to align model-based evaluations with actual human preferences.}
    \label{tab:criteria:problem}
    \vspace{-0.1in}
    \resizebox{0.95\textwidth}{!}{
    \renewcommand{\arraystretch}{1.1}
    \renewcommand{\tabcolsep}{2.5mm}
        \begin{tabular}{lll}
        \toprule
        \multicolumn{1}{p{.12\textwidth}}{\textbf{Types}} & \makecell{\multicolumn{1}{p{.12\textwidth}}{\textbf{Criteria}}} & \makecell{\multicolumn{1}{p{1.05\textwidth}}{\textbf{Texts}}} \\
        \midrule
        \multirow{39}{*}{\textbf{Problem}} 
        & \makecell{\multicolumn{1}{p{.12\textwidth}}{\textbf{Clarity}}} & \makecell{
            \multicolumn{1}{p{1.05\textwidth}}{1. The problem is presented in a highly ambiguous manner, lacking clear definition and leaving significant room for interpretation or confusion.} \\
            \multicolumn{1}{p{1.05\textwidth}}{2. The problem is somewhat defined but suffers from vague terms and insufficient detail, making it challenging to grasp the full scope or objective.} \\
            \multicolumn{1}{p{1.05\textwidth}}{3. The problem is stated in a straightforward manner, but lacks the depth or specificity needed to fully convey the nuances and boundaries of the research scope.} \\
            \multicolumn{1}{p{1.05\textwidth}}{4. The problem is clearly articulated with precise terminology and sufficient detail, providing a solid understanding of the scope and objectives with minimal ambiguity.} \\
            \multicolumn{1}{p{1.05\textwidth}}{5. The problem is exceptionally clear, concise, and specific, with every term and aspect well-defined, leaving no room for misinterpretation and fully encapsulating the research scope and aims.}
        } \\
        \noalign{\vskip 0.25ex}\cdashline{2-3}\noalign{\vskip 0.75ex}
        & \makecell{\multicolumn{1}{p{.12\textwidth}}{\textbf{Relevance}}} & \makecell{
            \multicolumn{1}{p{1.05\textwidth}}{1. The problem shows almost no relevance to the current field, failing to connect with the established context or build upon existing work.} \\
            \multicolumn{1}{p{1.05\textwidth}}{2. The problem has minimal relevance, with only superficial connections to the field and a lack of meaningful integration with prior studies.} \\
            \multicolumn{1}{p{1.05\textwidth}}{3. The problem is somewhat relevant, making a moderate attempt to align with the field but lacking significant innovation or depth.} \\
            \multicolumn{1}{p{1.05\textwidth}}{4. The problem is relevant and well-connected to the field, demonstrating a good understanding of existing work and offering promising contributions.} \\
            \multicolumn{1}{p{1.05\textwidth}}{5. The problem is highly relevant, deeply integrated with the current context, and represents a significant advancement in the field.}
        } \\
        \noalign{\vskip 0.25ex}\cdashline{2-3}\noalign{\vskip 0.75ex}
        & \makecell{\multicolumn{1}{p{.12\textwidth}}{\textbf{Originality}}} & \makecell{
            \multicolumn{1}{p{1.05\textwidth}}{1. The problem exhibits no discernible originality, closely mirroring existing studies without introducing any novel perspectives or challenges.} \\
            \multicolumn{1}{p{1.05\textwidth}}{2. The problem shows minimal originality, with slight variations from known studies, lacking significant new insights or innovative approaches.} \\
            \multicolumn{1}{p{1.05\textwidth}}{3. The problem demonstrates moderate originality, offering some new insights or angles, but these are not sufficiently groundbreaking or distinct from existing work.} \\
            \multicolumn{1}{p{1.05\textwidth}}{4. The problem is notably original, presenting a unique challenge or perspective that is well-differentiated from existing studies, contributing valuable new understanding to the field.} \\
            \multicolumn{1}{p{1.05\textwidth}}{5. The problem is highly original, introducing a pioneering challenge or perspective that has not been explored before, setting a new direction for future research.}
        } \\
        \noalign{\vskip 0.25ex}\cdashline{2-3}\noalign{\vskip 0.75ex}
        & \makecell{\multicolumn{1}{p{.12\textwidth}}{\textbf{Feasibility}}} & \makecell{
            \multicolumn{1}{p{1.05\textwidth}}{1. The problem is fundamentally infeasible due to insurmountable resource constraints, lack of foundational research, or critical methodological flaws.} \\
            \multicolumn{1}{p{1.05\textwidth}}{2. The problem faces significant feasibility challenges related to resource availability, existing knowledge gaps, or technical limitations, making progress unlikely.} \\
            \multicolumn{1}{p{1.05\textwidth}}{3. The problem is feasible to some extent but faces notable obstacles in resources, existing research support, or technical implementation, which could hinder significant advancements.} \\
            \multicolumn{1}{p{1.05\textwidth}}{4. The problem is mostly feasible with manageable challenges in resources, supported by adequate existing research, and has a clear, achievable methodology, though minor issues may persist.} \\
            \multicolumn{1}{p{1.05\textwidth}}{5. The problem is highly feasible with minimal barriers, well-supported by existing research, ample resources, and a robust, clear methodology, promising significant advancements.}
        } \\
        \noalign{\vskip 0.25ex}\cdashline{2-3}\noalign{\vskip 0.75ex}
        & \makecell{\multicolumn{1}{p{.12\textwidth}}{\textbf{Significance}}} & \makecell{
            \multicolumn{1}{p{1.05\textwidth}}{1. The problem shows minimal to no significance, lacking relevance or potential impact in advancing the field or contributing to practical applications.} \\
            \multicolumn{1}{p{1.05\textwidth}}{2. The problem has limited significance, with a narrow scope of impact and minor contributions to the field, offering little to no practical implications.} \\
            \multicolumn{1}{p{1.05\textwidth}}{3. The problem demonstrates average significance, with some contributions to the field and potential practical implications, but lacks innovation or broader impact.} \\
            \multicolumn{1}{p{1.05\textwidth}}{4. The problem is significant, offering notable contributions to the field and valuable practical implications, with evidence of potential for broader impact and advancement.} \\
            \multicolumn{1}{p{1.05\textwidth}}{5. The problem presents exceptional significance, with groundbreaking contributions to the field, broad and transformative potential impacts, and substantial practical applications across diverse domains.}
        } \\
        \bottomrule
        \end{tabular}
    }
\end{table*}

\begin{table*}
    \caption{The criteria induced from human judgments for validating the developed methods, which used to align model-based evaluations with actual human preferences.}
    \label{tab:criteria:method}
    \vspace{-0.1in}
    \resizebox{0.95\textwidth}{!}{
    \renewcommand{\arraystretch}{1.1}
    \renewcommand{\tabcolsep}{2.5mm}
        \begin{tabular}{lll}
        \toprule
        \multicolumn{1}{p{.12\textwidth}}{\textbf{Types}} & \makecell{\multicolumn{1}{p{.12\textwidth}}{\textbf{Criteria}}} & \makecell{\multicolumn{1}{p{1.05\textwidth}}{\textbf{Texts}}} \\
        \midrule
        \multirow{39}{*}{\textbf{Method}} 
        & \makecell{\multicolumn{1}{p{.12\textwidth}}{\textbf{Clarity}}} & \makecell{
            \multicolumn{1}{p{1.05\textwidth}}{1. The method is explained in an extremely vague or ambiguous manner, making it impossible to understand or replicate the approach without additional information or clarification.} \\
            \multicolumn{1}{p{1.05\textwidth}}{2. The method is described with some detail, but significant gaps in explanation or logic leave the reader with considerable confusion and uncertainty about how to apply or replicate the approach.} \\
            \multicolumn{1}{p{1.05\textwidth}}{3. The method is described with sufficient detail to understand the basic approach, but lacks the precision or specificity needed to fully replicate or grasp the nuances of the methodology without further guidance.} \\
            \multicolumn{1}{p{1.05\textwidth}}{4. The method is clearly and precisely described, with most details provided to allow for replication and comprehension, though minor areas may benefit from further clarification or elaboration.} \\
            \multicolumn{1}{p{1.05\textwidth}}{5. The method is articulated in an exceptionally clear, precise, and detailed manner, enabling straightforward replication and thorough understanding of the approach with no ambiguities.}
        } \\
        \noalign{\vskip 0.25ex}\cdashline{2-3}\noalign{\vskip 0.75ex}
        & \makecell{\multicolumn{1}{p{.12\textwidth}}{\textbf{Validity}}} & \makecell{
            \multicolumn{1}{p{1.05\textwidth}}{1. The method shows a fundamental misunderstanding of the research problem and lacks any credible alignment with established scientific principles or relevant studies.} \\
            \multicolumn{1}{p{1.05\textwidth}}{2. The method partially addresses the research problem but exhibits significant flaws in its scientific underpinning, making its validity questionable despite some alignment with existing literature.} \\
            \multicolumn{1}{p{1.05\textwidth}}{3. The method adequately addresses the research problem but with some limitations in its scientific validity, showing a mix of strengths and weaknesses in its alignment with related studies.} \\
            \multicolumn{1}{p{1.05\textwidth}}{4. The method effectively addresses the research problem, demonstrating a strong scientific basis and sound alignment with existing literature, albeit with minor areas for improvement.} \\
            \multicolumn{1}{p{1.05\textwidth}}{5. The method exemplifies an exceptional understanding of the research problem, grounded in a robust scientific foundation, and shows exemplary integration and advancement of existing studies' findings.}
        } \\
        \noalign{\vskip 0.25ex}\cdashline{2-3}\noalign{\vskip 0.75ex}
        & \makecell{\multicolumn{1}{p{.12\textwidth}}{\textbf{Rigorousness}}} & \makecell{
            \multicolumn{1}{p{1.05\textwidth}}{1. The method demonstrates a fundamental lack of systematic approach, with significant inconsistencies and inaccuracies in addressing the research problem, showing a disregard for established research standards.} \\
            \multicolumn{1}{p{1.05\textwidth}}{2. The method shows a minimal level of systematic effort but is marred by notable inaccuracies, lack of precision, and inconsistencies that undermine the rigorousness of the method in tackling the research problem.} \\
            \multicolumn{1}{p{1.05\textwidth}}{3. The method exhibits an average level of systematic structure and adherence to research standards but lacks the thoroughness, precision, and consistency required for a rigorous scientific inquiry.} \\
            \multicolumn{1}{p{1.05\textwidth}}{4. The method is well-structured and systematic, with a good level of precision and consistency, indicating a strong adherence to research standards, though it falls short of exemplifying the highest level of rigorousness.} \\
            \multicolumn{1}{p{1.05\textwidth}}{5. The method exemplifies exceptional rigorousness, with outstanding thoroughness, precision, and consistency in its systematic approach, setting a benchmark for high standards in scientific research quality.}
        } \\
        \noalign{\vskip 0.25ex}\cdashline{2-3}\noalign{\vskip 0.75ex}
        & \makecell{\multicolumn{1}{p{.12\textwidth}}{\textbf{Innovativeness}}} & \makecell{
            \multicolumn{1}{p{1.05\textwidth}}{1. The method introduces no novel elements, fully relying on existing techniques without any attempt to modify or adapt them for the specific research problem, showing a lack of innovativeness.} \\
            \multicolumn{1}{p{1.05\textwidth}}{2. The method shows minimal innovation, with only slight modifications to existing techniques that do not substantially change or improve the approach to the research problem.} \\
            \multicolumn{1}{p{1.05\textwidth}}{3. The method demonstrates moderate innovativeness, incorporating known techniques with some new elements or combinations that offer a somewhat fresh approach to the research problem but fall short of a significant breakthrough.} \\
            \multicolumn{1}{p{1.05\textwidth}}{4. The method is highly innovative, introducing new techniques or novel combinations of existing methods that significantly differ from standard practices, offering a new perspective or solution to the research problem.} \\
            \multicolumn{1}{p{1.05\textwidth}}{5. The method represents a groundbreaking innovation, fundamentally transforming the approach to the research problem with novel techniques or methodologies that redefine the field's standard practices.}
        } \\
        \noalign{\vskip 0.25ex}\cdashline{2-3}\noalign{\vskip 0.75ex}
        & \makecell{\multicolumn{1}{p{.12\textwidth}}{\textbf{Generalizability}}} & \makecell{
            \multicolumn{1}{p{1.05\textwidth}}{1. The method shows no adaptability, failing to extend its applicability beyond its original context or dataset, showing a complete lack of generalizability.} \\
            \multicolumn{1}{p{1.05\textwidth}}{2. The method demonstrates minimal adaptability, with limited evidence of potential applicability to contexts slightly different from the original.} \\
            \multicolumn{1}{p{1.05\textwidth}}{3. The method exhibits some level of adaptability, suggesting it could be applicable to related contexts or datasets with modifications.} \\
            \multicolumn{1}{p{1.05\textwidth}}{4. The method is adaptable and shows evidence of applicability to a variety of contexts or datasets beyond the original.} \\
            \multicolumn{1}{p{1.05\textwidth}}{5. The method is highly adaptable, demonstrating clear evidence of broad applicability across diverse contexts, populations, and settings.}
        } \\
        \bottomrule
        \end{tabular}
    }
\end{table*}

\begin{table*}
    \caption{The criteria induced from human judgments for validating the experiment designs, which are used to align model-based evaluations with actual human preferences.}
    \label{tab:criteria:experiment}
    \vspace{-0.1in}
    \resizebox{0.95\textwidth}{!}{
    \renewcommand{\arraystretch}{1.1}
    \renewcommand{\tabcolsep}{2.5mm}
        \begin{tabular}{lll}
        \toprule
        \multicolumn{1}{p{.12\textwidth}}{\textbf{Types}} & \makecell{\multicolumn{1}{p{.12\textwidth}}{\textbf{Criteria}}} & \makecell{\multicolumn{1}{p{1.05\textwidth}}{\textbf{Texts}}} \\
        \midrule
        \multirow{50}{*}{\textbf{Experiment}} 
        & \makecell{\multicolumn{1}{p{.12\textwidth}}{\textbf{Clarity}}} & \makecell{
            \multicolumn{1}{p{1.05\textwidth}}{1. The experiment design is extremely unclear, with critical details missing or ambiguous, making it nearly impossible for others to understand the setup, procedure, or expected outcomes.} \\
            \multicolumn{1}{p{1.05\textwidth}}{2. The experiment design lacks significant clarity, with many important aspects poorly explained or omitted, challenging others to grasp the essential elements of the setup, procedure, or expected outcomes.} \\
            \multicolumn{1}{p{1.05\textwidth}}{3. The experiment design is moderately clear, but some aspects are not detailed enough, leaving room for interpretation or confusion about the setup, procedure, or expected outcomes.} \\
            \multicolumn{1}{p{1.05\textwidth}}{4. The experiment design is mostly clear, with most aspects well-described, allowing others to understand the setup, procedure, and expected outcomes with minimal ambiguity.} \\
            \multicolumn{1}{p{1.05\textwidth}}{5. The experiment design is exceptionally clear, precise, and detailed, enabling easy understanding of the setup, procedure, and expected outcomes, with no ambiguity or need for further clarification.}
        } \\
        \noalign{\vskip 0.25ex}\cdashline{2-3}\noalign{\vskip 0.75ex}
        & \makecell{\multicolumn{1}{p{.12\textwidth}}{\textbf{Validity}}} & \makecell{
            \multicolumn{1}{p{1.05\textwidth}}{1. The experiment design demonstrates a fundamental misunderstanding of the research problem, lacks alignment with scientific methods, and shows no evidence of validity in addressing the research questions or testing the proposed methods.} \\
            \multicolumn{1}{p{1.05\textwidth}}{2. The experiment design has significant flaws in its approach to the research problem and scientific method, with minimal or questionable evidence of validity, making it largely ineffective in addressing the research questions or testing the proposed methods.} \\
            \multicolumn{1}{p{1.05\textwidth}}{3. The experiment design is generally aligned with the research problem and scientific method but has some limitations in its validity, offering moderate evidence that it can somewhat effectively address the research questions or test the proposed methods.} \\
            \multicolumn{1}{p{1.05\textwidth}}{4. The experiment design is well-aligned with the research problem and scientific method, providing strong evidence of validity and effectively addressing the research questions and testing the proposed methods, despite minor limitations.} \\
            \multicolumn{1}{p{1.05\textwidth}}{5. The experiment design excellently aligns with the research problem and scientific method, demonstrating robust evidence of validity and outstandingly addressing the research questions and testing the proposed methods without significant limitations.}
        } \\
        \noalign{\vskip 0.25ex}\cdashline{2-3}\noalign{\vskip 0.75ex}
        & \makecell{\multicolumn{1}{p{.12\textwidth}}{\textbf{Robustness}}} & \makecell{
            \multicolumn{1}{p{1.05\textwidth}}{1. The experiment design demonstrates a fundamental lack of understanding of the scientific method, with no evidence of durability or adaptability across varying conditions, leading to highly unreliable and non-replicable results.} \\
            \multicolumn{1}{p{1.05\textwidth}}{2. The experiment design shows minimal consideration for robustness, with significant oversights in addressing variability and ensuring consistency across different scenarios, resulting in largely unreliable outcomes.} \\
            \multicolumn{1}{p{1.05\textwidth}}{3. The experiment design adequately addresses some aspects of robustness but lacks comprehensive measures to ensure durability and consistency across a wide range of conditions, leading to moderate reliability.} \\
            \multicolumn{1}{p{1.05\textwidth}}{4. The experiment design incorporates a solid understanding of robustness, with clear efforts to ensure the experiment's durability and consistency across diverse conditions, though minor improvements are still possible for optimal reliability.} \\
            \multicolumn{1}{p{1.05\textwidth}}{5. The experiment design exemplifies an exceptional commitment to robustness, with meticulous attention to durability and adaptability across all possible conditions, ensuring highly reliable and universally applicable results.}
        } \\
        \noalign{\vskip 0.25ex}\cdashline{2-3}\noalign{\vskip 0.75ex}
        & \makecell{\multicolumn{1}{p{.12\textwidth}}{\textbf{Feasibility}}} & \makecell{
            \multicolumn{1}{p{1.05\textwidth}}{1. The experiment design is fundamentally unfeasible, with insurmountable resource, time, or technological constraints that make implementation virtually impossible within the proposed framework.} \\
            \multicolumn{1}{p{1.05\textwidth}}{2. The experiment design faces significant feasibility challenges, including major resource, time, or technological limitations, that heavily compromise its practical execution and likelihood of success.} \\
            \multicolumn{1}{p{1.05\textwidth}}{3. The experiment design is somewhat feasible, with moderate constraints on resources, time, or technology that could be addressed with adjustments, though these may not guarantee success.} \\
            \multicolumn{1}{p{1.05\textwidth}}{4. The experiment design is largely feasible, with minor resource, time, or technological limitations that can be effectively managed or mitigated, ensuring a high probability of successful implementation.} \\
            \multicolumn{1}{p{1.05\textwidth}}{5. The experiment design is highly feasible, with no significant constraints on resources, time, or technology, indicating that it can be implemented smoothly and successfully within the proposed framework.}
        } \\
        \noalign{\vskip 0.25ex}\cdashline{2-3}\noalign{\vskip 0.75ex}
        & \makecell{\multicolumn{1}{p{.12\textwidth}}{\textbf{Reproducibility}}} & \makecell{
            \multicolumn{1}{p{1.05\textwidth}}{1. The experiment design lacks critical details, making it virtually impossible for other researchers to replicate the study under the same conditions or methodologies.} \\
            \multicolumn{1}{p{1.05\textwidth}}{2. The experiment provides some essential information but omits significant details needed for replication, leading to considerable ambiguity in methodology or conditions.} \\
            \multicolumn{1}{p{1.05\textwidth}}{3. The experiment design includes sufficient details for replication, but lacks clarity or completeness in certain areas, posing challenges for seamless reproducibility.} \\
            \multicolumn{1}{p{1.05\textwidth}}{4. The experiment is well-documented with clear, detailed instructions and methodologies that allow for consistent replication, albeit with minor areas for improvement.} \\
            \multicolumn{1}{p{1.05\textwidth}}{5. The experiment design is exemplary in its clarity, detail, and comprehensiveness, ensuring that other researchers can precisely and effortlessly replicate the study under identical conditions and methodologies.}
        } \\
        \bottomrule
        \end{tabular}
    }
\end{table*}
\onecolumn

\begin{center}

\begingroup


\fontsize{6.5pt}{9pt}\selectfont

% \setlength\LTleft{-0.5in}
% \setlength\LTcapwidth{\textwidth} % default: 4in (rather less than \textwidth...)      
% \setlength\LTleft{0pt}            % default: \fill
% \setlength\LTright{0pt}           % default: \fill  

\begin{longtable}{p{.05\textwidth}p{.08\textwidth}p{.75\textwidth}}
\caption{The examples of research idea generation results from the proposed full ResearchAgent.} 
\vspace{-0.1in}
\label{tab:examples} \\

\hline 
\textbf{Index} & \textbf{Types} & \hspace{0.05in} \textbf{Texts}
\\ \hline 
\endfirsthead

\multicolumn{3}{c}%
{{\bfseries \tablename\ \thetable{} -- Continued from the previous page}} \\
\hline 
\textbf{Index} & \textbf{Types} & \hspace{0.05in} \textbf{Texts}
\\ \hline 
\endhead 

\multicolumn{3}{r}{\textit{\textbf{Continued on the next page}}} \\ \hline
\endfoot

\hline \hline
\endlastfoot

%%%%%%%%%%%%%%%%%%%%%%%%%%%%%% Examples %%%%%%%%%%%%%%%%%%%
\multirow{4}{*}{\textbf{1}} & 
\multicolumn{1}{p{.08\textwidth}}{\textbf{Input}} & 
\makecell{
    \multicolumn{1}{p{.75\textwidth}}{\textbf{Title}: Knowledge-Augmented Language Model Prompting for Zero-Shot Knowledge Graph Question Answering} \\
    \noalign{\vskip 0.25ex}\cdashline{0-0}\noalign{\vskip 0.75ex}
    \multicolumn{1}{p{.75\textwidth}}{\textbf{Abstract}: Large Language Models (LLMs) are capable of performing zero-shot closed-book question answering tasks, based on their internal knowledge stored in parameters during pre-training. However, such internalized knowledge might be insufficient and incorrect, which could lead LLMs to generate factually wrong answers. Furthermore, fine-tuning LLMs to update their knowledge is expensive. To this end, we propose to augment the knowledge directly in the input of LLMs. Specifically, we first retrieve the relevant facts to the input question from the knowledge graph based on semantic similarities between the question and its associated facts. After that, we prepend the retrieved facts to the input question in the form of the prompt, which is then forwarded to LLMs to generate the answer. Our framework, Knowledge-Augmented language model PromptING (KAPING), requires no model training, thus completely zero-shot. We validate the performance of our KAPING framework on the knowledge graph question answering task, that aims to answer the user’s question based on facts over a knowledge graph, on which ours outperforms relevant zero-shot baselines by up to 48\% in average, across multiple LLMs of various sizes.} \\
    \noalign{\vskip 0.25ex}\cdashline{0-0}\noalign{\vskip 0.75ex}
    \multicolumn{1}{p{.75\textwidth}}{\textbf{Entities}: ['Natural language', 'Learning management system', 'Question answering', 'Wikipedia', 'Artificial intelligence', 'Named-entity recognition', 'Quality assurance', 'Long short-term memory', 'English language', 'Google Books', 'Natural-language programming', 'LMS color space', 'GitHub', 'Logic learning machine', 'Machine translation', 'DBpedia', 'Integrated library system', 'Spanish language', 'Neural machine translation', 'ArXiv', 'Artificial general intelligence', 'Natural-language understanding', 'Nonlinear programming', 'Codex', 'LoRa', 'Adhesive', 'Common Crawl', 'ROUGE (metric)', 'Transformer', 'Tomas Mikolov']} \\
} \\
\cline{2-3}
& \multicolumn{1}{p{.08\textwidth}}{\textbf{Problem}} & 
\makecell{
    \multicolumn{1}{p{.75\textwidth}}{Developing a Multimodal Knowledge-Aware Prompting System for Multilingual Zero-Shot Question Answering across Structured and Unstructured Data Sources} 
} \\
\cline{2-3}
& \multicolumn{1}{p{.08\textwidth}}{\textbf{Method}} &
\makecell{
    \multicolumn{1}{p{.75\textwidth}}{The further refined and enhanced method for developing a multimodal knowledge-aware prompting system for multilingual zero-shot question answering (MMKP-ZSQA) incorporates advanced refinements in the following components to address the requirements for innovativeness, clarity, and generalizability:} \\ \\ 
    \multicolumn{1}{p{.75\textwidth}}{1. Interconnected Multimodal Knowledge Orchestration (IMKO): } \\ 
    \multicolumn{1}{p{.75\textwidth}}{- IMKO is an advanced evolution of SMKI, incorporating state-of-the-art language and vision models to handle data verbalization with context-aware disambiguation methods.} \\
    \multicolumn{1}{p{.75\textwidth}}{- Techniques like attention mechanisms are employed to weigh the relevance of different data pieces when integrated, revealing how textual and visual data augment each other in a synthesized representation.} \\ \\
    \multicolumn{1}{p{.75\textwidth}}{2. Linguistically Inclusive Retrieval Engine (LIRE): } \\ 
    \multicolumn{1}{p{.75\textwidth}}{- LIRE extends EMKA with an emphasis on semantic understanding, using transformer-based models trained on diverse datasets, including idiomatic and cultural nuances across languages.} \\
    \multicolumn{1}{p{.75\textwidth}}{- Specific algorithms to handle linguistic phenomena such as code-switching and transliteration are included, enhancing the application to a broader set of languages and dialects.} \\ \\
    \multicolumn{1}{p{.75\textwidth}}{3. Prompt Learning and Optimization Nexus (PLON): } \\ 
    \multicolumn{1}{p{.75\textwidth}}{- Building on AMPL, PLON focuses on creating a library of optimized prompts categorized by linguistic features and data modalities, using Bayesian optimization algorithms.} \\
    \multicolumn{1}{p{.75\textwidth}}{- It includes domain adaptation techniques and a wider array of meta-learning strategies with case studies for high and low-resource languages, and outlines theoretical frameworks for their implementation.} \\ \\ 
    \multicolumn{1}{p{.75\textwidth}}{4. Cross-Modal Integrative Learning System (C-MILS): } \\
    \multicolumn{1}{p{.75\textwidth}}{- C-MILS advances SCAS by detailing the use of multi-head attention across modalities for effective and scalable reasoning, with exemplifications on how each modality can enhance comprehension synergistically.} \\
    \multicolumn{1}{p{.75\textwidth}}{- The component now incorporates a layer of abstraction to distill knowledge into a modality-agnostic format, aiding reasoning and facilitating interpretability across languages and data types.} \\ \\ 
    \multicolumn{1}{p{.75\textwidth}}{5. User-Centric Adaptation and Privacy Framework (UCAPF): } \\
    \multicolumn{1}{p{.75\textwidth}}{- UCAPF enriches ALCUM with a user-centric design, focusing on interaction protocols that describe user engagement cycles, feedback loops, and privacy-preserving active learning paradigms, with practical workflows and pseudocode.} \\
    \multicolumn{1}{p{.75\textwidth}}{- Aligns with GDPR and other privacy frameworks to fortify trust and address the ethical use of data in dynamic learning environments.} \\ \\ 
    \multicolumn{1}{p{.75\textwidth}}{6. Global Evaluation and Reflective Testing System (GERTS): } \\
    \multicolumn{1}{p{.75\textwidth}}{- GERTS expands CEB by presenting a multi-tiered validation and reflection methodology to adjust system components based on a diverse set of metrics, including fairness, bias, interpretability, and computational efficiency.} \\
    \multicolumn{1}{p{.75\textwidth}}{- Offers a structured approach to cross-cultural evaluation, including the use of demographic and regional diversity in forming test cohorts.} \\ 
} \\
\cline{2-3}
& \multicolumn{1}{p{.08\textwidth}}{\textbf{Experiment}} & 
\makecell{
    \multicolumn{1}{p{.75\textwidth}}{The experiment, named "Refined Experiment for Multimodal Knowledge-Aware Prompting System for Multilingual Zero-shot Question Answering (RE-MKP-ZSQA)", aims to methodically develop and validate an advanced AI system. The experiment is streamlined to address feasibility, clarity, and reproducibility concerns while upholding robustness and validity by adhering to the following refined phases:} \\ \\ 
    \multicolumn{1}{p{.75\textwidth}}{1. Detailed System Implementation Plan:} \\
    \multicolumn{1}{p{.75\textwidth}}{- Provide a publicly accessible project roadmap with specific milestones, resource allocation, and timelines.} \\ \\
    \multicolumn{1}{p{.75\textwidth}}{2. Dataset Curation with Clear Guidelines:} \\
    \multicolumn{1}{p{.75\textwidth}}{- Publish precise annotation guidelines with strategies to prevent bias.} \\
    \multicolumn{1}{p{.75\textwidth}}{- Document the dataset assembly process, including source selection and data processing procedures.} \\ \\
    \multicolumn{1}{p{.75\textwidth}}{3. Transparent System Training:} \\
    \multicolumn{1}{p{.75\textwidth}}{- Offer a detailed training protocol with hyperparameters, optimization strategies, and Bayesian optimization processes used in PLON.} \\ \\
    \multicolumn{1}{p{.75\textwidth}}{4. Structured Zero-Shot Evaluation:} \\
    \multicolumn{1}{p{.75\textwidth}}{- Outline evaluation metrics derived from GERTS with benchmark datasets to test zero-shot capabilities.} \\ \\
    \multicolumn{1}{p{.75\textwidth}}{5. Clearer Interdisciplinary Evaluation Protocol:} \\
    \multicolumn{1}{p{.75\textwidth}}{- Specify the composition of the evaluation committee, criteria for assessments, and methods for integrating the feedback.} \\ \\
    \multicolumn{1}{p{.75\textwidth}}{6. Iterative Improvement with Validation Metrics:} \\
    \multicolumn{1}{p{.75\textwidth}}{- Describe statistical methods for reflective assessment and continuous improvement, aligned with multi-tiered GERTS methodology.} \\ \\
    \multicolumn{1}{p{.75\textwidth}}{7. User-Centric Design and Privacy Compliance Evaluation:} \\
    \multicolumn{1}{p{.75\textwidth}}{- Structure user studies with targeted data points to assess usability and cultural adaptability.} \\
    \multicolumn{1}{p{.75\textwidth}}{- Outline privacy compliance protocols to adhere to international standards.} \\ \\
    \multicolumn{1}{p{.75\textwidth}}{8. Detailed Global Scalability Evaluation Method:} \\
    \multicolumn{1}{p{.75\textwidth}}{- Define evaluation metrics for scalability tests and describe diverse infrastructural setups.} \\ \\
    \multicolumn{1}{p{.75\textwidth}}{9. Enhanced Reporting for Reproducibility:} \\
    \multicolumn{1}{p{.75\textwidth}}{- Commit to creating a comprehensive report with precise specifications, configurations, and instructions for replication purposes.} \\
    \multicolumn{1}{p{.75\textwidth}}{- Utilize GitHub for version-controlled deposition of code and datasets, and arXiv for openly accessible experiment protocols and findings.} \\
}
\\

\midrule

\multirow{25}{*}{\textbf{2}} & 
\multicolumn{1}{p{.08\textwidth}}{\textbf{Input}} & 
\makecell{
    \multicolumn{1}{p{.75\textwidth}}{\textbf{Title}: Test-Time Self-Adaptive Small Language Models for Question Answering} \\
    \noalign{\vskip 0.25ex}\cdashline{0-0}\noalign{\vskip 0.75ex}
    \multicolumn{1}{p{.75\textwidth}}{\textbf{Abstract}: Recent instruction-finetuned large language models (LMs) have achieved notable performances in various tasks, such as question-answering (QA). However, despite their ability to memorize a vast amount of general knowledge across diverse tasks, they might be suboptimal on specific tasks due to their limited capacity to transfer and adapt knowledge to target tasks. Moreover, further finetuning LMs with labeled datasets is often infeasible due to their absence, but it is also questionable if we can transfer smaller LMs having limited knowledge only with unlabeled test data. In this work, we show and investigate the capabilities of smaller self-adaptive LMs, only with unlabeled test data. In particular, we first stochastically generate multiple answers, and then ensemble them while filtering out low-quality samples to mitigate noise from inaccurate labels. Our proposed self-adaption strategy demonstrates significant performance improvements on benchmark QA datasets with higher robustness across diverse prompts, enabling LMs to stay stable.} \\
    \noalign{\vskip 0.25ex}\cdashline{0-0}\noalign{\vskip 0.75ex}
    \multicolumn{1}{p{.75\textwidth}}{\textbf{Entities}: ['Codex', 'Natural language', 'English language', 'United States', 'Question answering', 'Natural-language programming', 'GTRI Information and Communications Laboratory', 'Artificial intelligence', 'LoRa', 'Llama', 'Python (programming language)', 'Learning management system', 'Natural language processing', 'Reinforcement learning', 'LMS color space', 'Wikipedia', 'GitHub', 'Natural-language understanding', 'London, Midland and Scottish Railway', 'Integrated library system', 'Language model', 'Chinese language', 'Lumen (unit)', 'Spanish language', 'English Wikipedia', 'Logic learning machine', 'Gradient descent', 'Alternative public offering', 'Technology transfer', 'Dialogue system']} \\
} \\
\cline{2-3}
& \multicolumn{1}{p{.08\textwidth}}{\textbf{Problem}} & 
\makecell{
    \multicolumn{1}{p{.75\textwidth}}{Developing a Scalable, Domain-Adaptive Test-Time Training Protocol for Low-Resource Language QA Using Small Language Models} \\ 
} \\
\cline{2-3}
& \multicolumn{1}{p{.08\textwidth}}{\textbf{Method}} &
\makecell{
    \multicolumn{1}{p{.75\textwidth}}{1. Selection of Scalable Compact Language Models (CLMs): Identify and evaluate existing CLMs suitable for adaptation, emphasizing models with minimal computational requirements.} \\ \\
    \multicolumn{1}{p{.75\textwidth}}{2. Creation of a Multilingual Test-Time Training (TTT) Framework: Develop a TTT protocol that enables CLMs to adapt to new domains and languages during the inference phase, leveraging unsupervised learning techniques and pseudo-label generation.} \\ \\
    \multicolumn{1}{p{.75\textwidth}}{3. Synthetic and Unsupervised Data Generation: Utilize a combination of unsupervised and synthetic data generation methods to produce multilingual QA pairs, employing techniques such as back-translation and context-based question synthesis.} \\ \\
    \multicolumn{1}{p{.75\textwidth}}{4. Domain-Adaptive Mechanisms: Introduce domain-adaptive components, including feature adaptation layers and meta-learning algorithms, which tailor the model's behavior to new contexts and languages at test time.} \\ \\
    \multicolumn{1}{p{.75\textwidth}}{5. Incremental Language Addition and Dominance Assessment: Start with a subset of linguistically diverse, low-resource languages. Evaluate domain adaptability for each language via an iterative process, ensuring models learn to prioritize resource efficiency.} \\ \\
    \multicolumn{1}{p{.75\textwidth}}{6. Model Robustness and Generalization: Perform robustness tuning (RT) to prepare models for unforeseen linguistic variations and conduct thorough evaluations across multiple domains to ensure models can generalize their learning effectively.} \\ \\
    \multicolumn{1}{p{.75\textwidth}}{7. Human-In-The-Loop Evaluation: Conduct evaluations with native speakers and domain experts to validate the relevance and accuracy of the QA outputs, incorporating feedback into the iterative training process.} \\ \\
    \multicolumn{1}{p{.75\textwidth}}{8. Open-Sourcing and Community Collaboration: Make the TTT protocol, trained models, and evaluation benchmarks publicly available for the research community, fostering collaboration and further innovation.} \\ 
} \\
\cline{2-3}
& \multicolumn{1}{p{.08\textwidth}}{\textbf{Experiment}} & 
\makecell{ 
    \multicolumn{1}{p{.75\textwidth}}{1. Selection and Preparation:} \\
    \multicolumn{1}{p{.75\textwidth}}{- Identify potential compact language models (CLMs) suitable for domain adaptation and test-time training, focusing on those with minimal computational requirement and the ability to be fine-tuned or adapted in an unsupervised manner.} \\
    \multicolumn{1}{p{.75\textwidth}}{- Prepare a diverse set of low-resource languages and corresponding text corpora, ensuring linguistic diversity and sociocultural significance. Select benchmark datasets for these languages if available.} \\ \\
    \multicolumn{1}{p{.75\textwidth}}{2. Training and Adaptation Procedure:} \\
    \multicolumn{1}{p{.75\textwidth}}{- Create a Test-Time Training (TTT) framework that allows selected CLMs to adapt to various domains in the selected low-resource languages during the inference phase.} \\
    \multicolumn{1}{p{.75\textwidth}}{- Implement unsupervised learning techniques and pseudo-label generation to produce QA pairs, utilizing back-translation and context-based question synthesis to generate synthetic datasets for languages with limited or no available QA datasets.} \\
    \multicolumn{1}{p{.75\textwidth}}{- Integrate domain-adaptive components and meta-learning algorithms into the CLMs to enable domain-specific adaptations at test time.
} \\ \\
    \multicolumn{1}{p{.75\textwidth}}{3. Iterative Evaluation and Refinement:} \\
    \multicolumn{1}{p{.75\textwidth}}{- Begin adaptation and training with a single low-resource language and gradually add additional languages, monitoring the domain adaptability and model performance metrics after each addition.} \\ 
    \multicolumn{1}{p{.75\textwidth}}{- Perform robustness tuning and cross-domain evaluations for each CLM and language adaptation to ensure generalizability and prevent overfitting.} \\ \\
    \multicolumn{1}{p{.75\textwidth}}{4. Human-In-The-Loop Assessment:} \\
    \multicolumn{1}{p{.75\textwidth}}{- Enlist native speakers and domain experts to evaluate the relevance and accuracy of the model's QA outputs for each language.} \\ 
    \multicolumn{1}{p{.75\textwidth}}{- Incorporate feedback into the iterative training process, refining and re-adapting the models accordingly.} \\ \\
    \multicolumn{1}{p{.75\textwidth}}{5. Open-Sourcing and Community Feedback:} \\
    \multicolumn{1}{p{.75\textwidth}}{- Make the TTT protocol, adaptive CLMs, evaluation benchmarks, and any synthetic datasets publicly available for the research community.} \\ \\
    \multicolumn{1}{p{.75\textwidth}}{6. Experiment Monitoring and Documentation:} \\
    \multicolumn{1}{p{.75\textwidth}}{- Record all the parameters, datasets, model configurations, and evaluation metrics meticulously to ensure robustness and reproducibility.} \\ 
    \multicolumn{1}{p{.75\textwidth}}{- Document any challenges faced, unexpected results, or adaptions made during the experiment for open-sourcing purposes.} \\ \\
    \multicolumn{1}{p{.75\textwidth}}{7. Data Analysis and Reporting:} \\
    \multicolumn{1}{p{.75\textwidth}}{- Analyze the collected performance data quantitatively, using appropriate statistical methods to compare with non-adaptive baselines.} \\
    \multicolumn{1}{p{.75\textwidth}}{- Report qualitative findings from human-in-the-loop evaluations, interpreting the implications for language model performance in low-resource language domains.}
}
\\

\midrule

\multirow{2}{*}{\textbf{3}} & 
\multicolumn{1}{p{.08\textwidth}}{\textbf{Input}} & 
\makecell{
    \multicolumn{1}{p{.75\textwidth}}{\textbf{Title}: Whole-brain annotation and multi-connectome cell typing quantifies circuit stereotypy in Drosophila} \\
    \noalign{\vskip 0.25ex}\cdashline{0-0}\noalign{\vskip 0.75ex}
    \multicolumn{1}{p{.75\textwidth}}{\textbf{Abstract}: The fruit fly Drosophila melanogaster combines surprisingly sophisticated behaviour with a highly tractable nervous system. A large part of the fly’s success as a model organism in modern neuroscience stems from the concentration of collaboratively generated molecular genetic and digital resources. As presented in our FlyWire companion paper1, this now includes the first full brain connectome of an adult animal. Here we report the systematic and hierarchical annotation of this 130,000-neuron connectome including neuronal classes, cell types and developmental units (hemilineages). This enables any researcher to navigate this huge dataset and find systems and neurons of interest, linked to the literature through the Virtual Fly Brain database2. Crucially, this resource includes 4,552 cell types. 3,094 are rigorous consensus validations of cell types previously proposed in the “hemibrain” connectome3. In addition, we propose 1,458 new cell types, arising mostly from the fact that the FlyWire connectome spans the whole brain, whereas the hemibrain derives from a subvolume. Comparison of FlyWire and the hemibrain showed that cell type counts and strong connections were largely stable, but connection weights were surprisingly variable within and across animals. Further analysis defined simple heuristics for connectome interpretation: connections stronger than 10 unitary synapses or providing >1\% of the input to a target cell are highly conserved. Some cell types showed increased variability across connectomes: the most common cell type in the mushroom body, required for learning and memory, is almost twice as numerous in FlyWire as the hemibrain. We find evidence for functional homeostasis through adjustments of the absolute amount of excitatory input while maintaining the excitation-inhibition ratio. Finally, and surprisingly, about one third of the cell types proposed in the hemibrain connectome could not yet be reliably identified in the FlyWire connectome. We therefore suggest that cell types should be defined to be robust to inter-individual variation, namely as groups of cells that are quantitatively more similar to cells in a different brain than to any other cell in the same brain. Joint analysis of the FlyWire and hemibrain connectomes demonstrates the viability and utility of this new definition. Our work defines a consensus cell type atlas for the fly brain and provides both an intellectual framework and open source toolchain for brain-scale comparative connectomics.} \\
    \noalign{\vskip 0.25ex}\cdashline{0-0}\noalign{\vskip 0.75ex}
    \multicolumn{1}{p{.75\textwidth}}{\textbf{Entities}: ['Virtual Fly Brain', 'Central nervous system', 'Transposable element', 'SUMO protein', 'Kenyon cell', 'Romani people', 'Induced stem cells', 'Ventral nerve cord', 'FlyBase', "Parkinson's disease", 'Virtual Network Computing', 'P element', 'Piwi-interacting RNA', 'Drosophila Genetic Reference Panel', 'Bateson–Dobzhansky–Muller model', 'J. B. S. Haldane', 'ATG7', "Haldane's rule", 'Oxford Nanopore Technologies', 'Drosophila mauritiana', 'Germline', 'PINK1', 'Migratory locust', 'CRISPR', 'Heliconius', 'GINS (protein complex)', 'Parkin (ligase)', 'Lepidoptera', 'Illumina, Inc.', 'Drosophila']} \\
} \\
\cline{2-3}
& \multicolumn{1}{p{.08\textwidth}}{\textbf{Problem}} & 
\makecell{
    \multicolumn{1}{p{.75\textwidth}}{Investigating the Functional Implications of Connectome Variability in Drosophila's Learning and Memory Circuits Across Different Environmental and Genetic Contexts} \\ 
} \\
\cline{2-3}
& \multicolumn{1}{p{.08\textwidth}}{\textbf{Method}} &
\makecell{
    \multicolumn{1}{p{.75\textwidth}}{The proposed method involves a multi-tiered approach that integrates connectomics, behavioral assays, genetic manipulation, and computational modeling to investigate the functional implications of connectome variability in Drosophila's learning and memory circuits. The method consists of the following steps:} \\ \\
    \multicolumn{1}{p{.75\textwidth}}{1. Connectome Mapping and Variability Analysis:} \\ 
    \multicolumn{1}{p{.75\textwidth}}{a. Utilize the Virtual Fly Brain database to identify and compare individual connectomes, focusing on the mushroom body.} \\ 
    \multicolumn{1}{p{.75\textwidth}}{b. Quantify the variability in connection weights and cell type counts using statistical methods and machine learning algorithms to identify patterns of variability.} \\ \\
    \multicolumn{1}{p{.75\textwidth}}{2. Behavioral Assays:} \\
    \multicolumn{1}{p{.75\textwidth}}{a. Design a series of learning and memory tasks for Drosophila, such as olfactory conditioning or visual pattern recognition.} \\
    \multicolumn{1}{p{.75\textwidth}}{b. Test groups of flies with known connectome profiles under controlled environmental conditions to establish baseline behavioral data.} \\ \\
    \multicolumn{1}{p{.75\textwidth}}{3. Environmental and Genetic Perturbations:} \\ 
    \multicolumn{1}{p{.75\textwidth}}{a. Expose different groups of flies to varied learning paradigms and sensory inputs to create environmental perturbations.} \\ 
    \multicolumn{1}{p{.75\textwidth}}{b. Use CRISPR-Cas9 technology to introduce targeted mutations in genes like PINK1 or Parkin, creating genetic perturbations.} \\ 
    \multicolumn{1}{p{.75\textwidth}}{c. Assess the impact of these perturbations on connectome structure using high-resolution imaging and reconstruction techniques.} \\ \\
    \multicolumn{1}{p{.75\textwidth}}{4. Transcriptomic and Spatial Analysis:} \\ 
    \multicolumn{1}{p{.75\textwidth}}{a. Apply single-cell RNA sequencing and spatial transcriptomics to profile gene expression changes in response to environmental and genetic perturbations.} \\ 
    \multicolumn{1}{p{.75\textwidth}}{b. Correlate transcriptomic data with connectome changes to identify molecular pathways associated with structural and functional plasticity.} \\ \\
    \multicolumn{1}{p{.75\textwidth}}{5. Computational Modeling and Network Analysis:} \\ 
    \multicolumn{1}{p{.75\textwidth}}{a. Develop computational models to simulate the effects of connectome variability on neural circuit function.} \\ 
    \multicolumn{1}{p{.75\textwidth}}{b. Use network analysis tools to explore information flow and circuit dynamics, integrating data from related papers on olfactory projection neurons and information flow in the olfactory system.} \\ \\
    \multicolumn{1}{p{.75\textwidth}}{6. Synthesis and Validation:} \\ 
    \multicolumn{1}{p{.75\textwidth}}{a. Integrate findings from behavioral assays, connectome mapping, transcriptomic analysis, and computational modeling.} \\ 
    \multicolumn{1}{p{.75\textwidth}}{b. Validate the proposed models and hypotheses through iterative experimentation and refinement.} \\ 
} \\
\cline{2-3}
& \multicolumn{1}{p{.08\textwidth}}{\textbf{Experiment}} & 
\makecell{ 
    \multicolumn{1}{p{.75\textwidth}}{1. Pre-experimental Setup:} \\
    \multicolumn{1}{p{.75\textwidth}}{a. Establish a Drosophila melanogaster breeding program to ensure a consistent supply of genetically similar flies for the experiments.} \\
    \multicolumn{1}{p{.75\textwidth}}{b. Develop a standardized protocol for rearing flies under controlled conditions to minimize pre-experimental variability.} \\ 
    \multicolumn{1}{p{.75\textwidth}}{c. Select and validate CRISPR-Cas9 constructs for targeted gene editing of PINK1, Parkin, and other genes of interest.} \\ 
    \multicolumn{1}{p{.75\textwidth}}{d. Train personnel in the use of the Virtual Fly Brain database and relevant computational tools for connectome analysis.} \\ \\
    \multicolumn{1}{p{.75\textwidth}}{2. Connectome Mapping and Variability Analysis:} \\
    \multicolumn{1}{p{.75\textwidth}}{a. Randomly assign individual flies to either a control group or various treatment groups (environmental and genetic perturbations).} \\
    \multicolumn{1}{p{.75\textwidth}}{b. Utilize high-resolution imaging techniques to map the connectomes of flies from each group, with a focus on the mushroom body.} \\
    \multicolumn{1}{p{.75\textwidth}}{c. Apply statistical and machine learning algorithms to quantify and compare the variability in connection weights and cell type counts across groups.} \\ \\
    \multicolumn{1}{p{.75\textwidth}}{3. Behavioral Assays:} \\
    \multicolumn{1}{p{.75\textwidth}}{a. Design and validate a series of learning and memory tasks, such as olfactory conditioning and visual pattern recognition, ensuring tasks are sensitive to subtle differences in performance.} \\ 
    \multicolumn{1}{p{.75\textwidth}}{b. Test flies from each group in the behavioral tasks and record performance metrics.} \\ 
    \multicolumn{1}{p{.75\textwidth}}{c. Analyze behavioral data to establish correlations with connectome profiles.} \\ \\
    \multicolumn{1}{p{.75\textwidth}}{4. Environmental and Genetic Perturbations:} \\
    \multicolumn{1}{p{.75\textwidth}}{a. Expose flies to different learning paradigms and sensory inputs to induce environmental perturbations.} \\ 
    \multicolumn{1}{p{.75\textwidth}}{b. Perform gene editing using CRISPR-Cas9 to create genetic perturbations in the treatment groups.} \\ 
    \multicolumn{1}{p{.75\textwidth}}{c. Re-map connectomes post-perturbation to assess structural changes.} \\ \\
    \multicolumn{1}{p{.75\textwidth}}{5. Transcriptomic and Spatial Analysis:} \\
    \multicolumn{1}{p{.75\textwidth}}{a. Collect brain tissue from flies post-behavioral assays and perform single-cell RNA sequencing and spatial transcriptomics.} \\ 
    \multicolumn{1}{p{.75\textwidth}}{b. Analyze transcriptomic data to identify gene expression changes and correlate these with observed connectome and behavioral variations.} \\ \\
    \multicolumn{1}{p{.75\textwidth}}{6. Computational Modeling and Network Analysis:} \\
    \multicolumn{1}{p{.75\textwidth}}{a. Develop computational models to simulate the impact of observed connectome variability on neural circuit function.} \\ 
    \multicolumn{1}{p{.75\textwidth}}{b. Use network analysis to integrate behavioral, connectomic, and transcriptomic data, focusing on information flow and circuit dynamics.} \\ \\
    \multicolumn{1}{p{.75\textwidth}}{7. Synthesis and Validation:} \\
    \multicolumn{1}{p{.75\textwidth}}{a. Integrate findings across all experimental components to formulate a cohesive understanding of the functional implications of connectome variability.} \\
    \multicolumn{1}{p{.75\textwidth}}{b. Validate models and refine hypotheses through additional targeted experiments, informed by initial findings.}
}
\\

%%%%%%%%%%%%%%%%%%%%%%%%%%%%%%%%%%%%%%%%%%%%%%%%%%%%%%%%%%


\end{longtable}


\endgroup

\end{center}

\twocolumn
